\documentclass[12pt]{article}

\usepackage{amsmath}
\usepackage[margin=1in]{geometry}
\usepackage{fancyhdr}
\usepackage{amsthm}
\usepackage{braket}
\usepackage{amssymb}
\usepackage{tikz}
\usetikzlibrary{shapes.geometric}
\usepackage{enumerate}
\usepackage[shortlabels]{enumitem}


\pagestyle{fancy}
\fancyhead[l]{Hudson Benites}
\fancyhead[c]{Homework \#5}
\fancyfoot[c]{\thepage}
\renewcommand{\headrulewidth}{0.2pt}
\setlength{\headheight}{15pt}
\fancyhead[r]{\today}

\begin{document}
\begin{enumerate}[start = 1, label = {\bfseries Question \arabic*:}, leftmargin=1in]
       \item todo
       \item For the following proofs, let $R$ be a commutative ring and $x,y,z\in R$.\begin{itemize}
        \item[(a)] \begin{proof} We show that $(-x)(-y) + (-(xy)) = 0$, and thus $(-x)(-y)=xy$. As was shown in Proposition 5.7 in the notes, $-(xy) = (-x)y$. So, we compute that \begin{center}
            $(-x)(-y) + (-(xy)) = (-x)(-y) + (-x)y = (-x)(-y+y)=(-x)\cdot 0 = 0$.
        \end{center} Thus, we are done. \end{proof}

        \item[(b)] \begin{proof}
            Considering $x(y-z) = x(y+(-z))$, we can appeal to the distributive property of commutative rings. So, $x(y+(-z))=xy+x(-z)$. Appealing to Proposition 5.7 in the notes, we have that $x(-z)=-(xz)$, we get that $xy+x(-z)=xy-xz$.
            \end{proof}
        \item[(c)] \begin{proof}
            Considering $-1(x)$, we can appeal to Proposition 5.7 in the notes, and rewrite the expression as $-(1x)$. Appealing to $R$'s multiplicative identity, which is $1$, we know $1x = x, \forall x\in R$. Thus, $-(1x)=-x$.
        \end{proof}
       \end{itemize}    
       \item \begin{proof}
        Suppose $R$ is an ordered ring, $x,y,z \in R, x < y, 0 < z$. It follows that $y-x\in R^+$ and $z\in R^+$. Appealing to the closure of positives, $(y-x)(z)\in R^+$. Distributing, we get that $yz-xz\in R^+$. Thus, $xz<yz$.
       \end{proof}  
\end{enumerate}
\end{document}